\section{Reproductibilité, traçabilité et livrables}
\label{sec:reproduction}
\subsection{FAST : vérifier les artefacts de l'étude historique}
Le notebook Jupyter de Phase 9A a exécuté 52 cellules de revue des sorties gelées,
dans la référence et une copie sans données ni caches de modèles. Il affiche les
résultats et contrôle arithmétiquement les métriques depuis les prévisions
officielles. Cela ne reproduit pas les acquisitions, les estimations R/Python ni
les entraînements neuronaux. Les prévisions et recommandations des chapitres
précédents proviennent exclusivement des fichiers historiques.

\subsection{FULL : pipeline exécuté dans un sandbox préparé}
La préparation 9C1 a implémenté un notebook FULL et un orchestrateur séparés,
avec défaut PREFLIGHT et autorisation FULL explicite. Le run ultérieur a
réellement exécuté les scripts scientifiques dans le clone indépendant
\sourcepath{C:/Temp/ai-chip-full-9c1-e3cd311}, basé sur le point scientifique
\sourcehash{e3cd311e4177807c2d422e74f3290d93cef35dab}.
L'audit 9C3 relève \texttt{PIPELINE\_EXECUTED},
\texttt{PARTIAL\_PENDING\_CLEAN\_MACHINE\_REVIEW} et l'absence de fichier d'échec.

Les 17 subprocess scientifiques couvrent préparation, EDA, ARIMA et validation
R/Python, VAR/Johansen, GARCH, réseaux puis Phase 7. Tous leurs codes de sortie
sont zéro, ainsi que ceux du checkpoint et du gate. Le contrôle LLM interne est
documentaire : réponse intacte et 35 affirmations historiques présentes avec
sources existantes ; aucune nouvelle génération.
Les sept snapshots ont leurs SHA historiques. Les 358 fichiers protégés par le
manifeste de ce run sont inchangés. Cela décrit une preuve dans un environnement
préparé, pas une installation sur une machine vierge.

\begin{table}[!htbp]
\centering
\caption{Trois preuves de reproduction à distinguer au terme de 9C3.}
\label{tab:reproduction_status}
\small
\begin{tabularx}{\linewidth}{lXl}
\toprule
Preuve & État effectivement démontré & Statut \\
\midrule
A : sandbox & Runner scientifique FULL terminé, logs et résultats nouveaux & Exécuté \\
B : machine vierge & Packages Python/R copiés d'installations existantes & PARTIAL \\
C : notebook FULL & Seule une exécution PREFLIGHT du notebook est conservée & PARTIAL \\
§8.3 intégral & Notebook complet sur machine vierge non démontré & PARTIAL \\
\bottomrule
\end{tabularx}
\sourcecaption{Source : \sourcepath{docs/phase9c_full_execution_audit.md}.
Exécution du runner et identité numérique sont deux exigences différentes.}
\end{table}

\subsection{Différences observées, sans remplacement des résultats officiels}
Sur 210 fichiers, 116 sont identiques en bytes et 94 diffèrent : 64 fins de ligne,
20 durées/provenance, quatre images, un journal de réestimation et cinq fichiers
avec valeurs scientifiques différentes. Les contenus ont été caractérisés ;
le mécanisme numérique des variations GRU reste non établi.

Les refits Phase 6 conservent la même charge utile de poids HDF5 que les
historiques, mais les fichiers \texttt{.keras} ont des SHA différents par
leurs métadonnées/configurations sérialisées. Le loader original Phase 7
exigeant les empreintes de fichiers pré-déclarées, il a reconstruit les réseaux
C sur TRAIN avant l'accès TEST. Les poids numériques LSTM reconstruits sont
identiques ; 12865 paramètres GRU diffèrent, au maximum de
\(2{,}8312206268310547\times10^{-7}\).

Le nouveau panel modifie 304 prévisions GRU. Quatre écarts MSE/RMSE dépassent
la tolérance descriptive historique \(10^{-6}\), maintenue ; les quatre écarts
MAE/MAPE sont plus petits que ce seuil mais non nuls. Certaines valeurs DM et
ratios varient aussi. Les décisions DM/Holm, l'éligibilité et Naive restent
inchangés. Ces résultats de reproduction ne remplacent aucune valeur
officielle des tableaux du chapitre~\ref{sec:evaluation}.

La seed 42, les opérations déterministes demandées et TensorFlow 2.21.0 CPU ne
constituent pas une garantie d'identité entre environnements~\cite{R16}.
Le message oneDNN des logs ne prouve pas la cause de cette variation.
Ni tolérance, ni spécification, ni réponse LLM n'ont été modifiées pour obtenir
une correspondance.

\subsection{Prérequis et actions restantes pour §8.3}
Un essai futur sur machine vierge exige les snapshots originaux vérifiés, ou
un téléchargement contrôlé arrêté si ses bytes divergent ; Python 3.12.3,
les packages listés et leurs versions documentées ; R 4.5.2 avec forecast,
tseries, FinTS, vars, urca et rugarch ; TensorFlow CPU et les outils notebook.
Les scripts doivent être les sources réelles, avec les synthèses après les
exécutions R/Python. Le placeholder R EDA n'est pas une implémentation.

L'essai doit rester dans une nouvelle copie indépendante, avec l'historique et
les protocoles originaux, archivage des seules sorties copiées et espace Phase 7
initialement vide. Le checkpoint TRAIN précède le gate authentique, puis le
gel des modèles et l'ouverture TEST du nouvel évaluateur. L'ancien runtime
n'est pas une preuve de ce nouveau run. La réponse LLM officielle est chargée,
préservée et auditée comme document historique.

Le notebook doit terminer toutes ses cellules en mode FULL et conserver outputs,
logs, packages, hashes et provenance. Son paramètre \texttt{SOURCE\_REVISION}
doit correspondre au manifeste du nouveau clone ; l'exécuteur PREFLIGHT
n'injecte pas ce paramètre. Ces étapes sur machine vierge et notebook intégral
ne sont pas démontrées par l'exécution CLI existante.

Les 17 commandes du run audité cumulent 839,065 s, contre 845,982 s avec le
checkpoint/gate. Ces sommes excluent une partie des copies et contrôles et ne
sont pas le temps mural FULL. Le coût neuronal historique de 20,14 minutes est
un autre run et une autre mesure, pas une promesse de durée.

\subsection{Périmètre documentaire et version LaTeX}
Le PDF officiel du cahier des charges n'était pas intégré au dépôt lors de
l'audit automatisé. Cette limite ne signifie ni inexistence du cahier des
charges ni livrable scientifique obligatoire manquant. La checklist 9A reste
une preuve datée, avec un rapport alors différé ; le rapport Word/PDF 9B existe
désormais. Aucun total historique n'est présenté comme une certification
exhaustive actuelle des clauses du PDF.

Cette version LaTeX est additive : dix-sept images historiques copiées à SHA
identique, tableaux affichés depuis les CSV gelés, schéma TikZ documentaire et
annexe LLM fidèle. Elle ne modifie ni les exports Word/PDF de référence, ni
les protocoles, ni la recommandation. Le contenu est réparti en fichiers UTF-8
\texttt{.tex}, figures locales et BibTeX. L'utilisateur a importé et compilé
la version initiale dans Prism AI (29 pages affichées). La présente révision
est vérifiée localement ; sa compilation et sa relecture dans Prism restent à faire.

\sourcecaption{Sources : \sourcepath{docs/phase9a_reproducibility.md},
\sourcepath{docs/phase9c_preflight_audit.md},
\sourcepath{docs/phase9c_full_reproduction.md},
\sourcepath{docs/phase9c_full_execution_audit.md},
\sourcepath{docs/phase9c_full_summary.json},
\sourcepath{docs/final_requirements_checklist.md}. Les audits 9A/9B/9C1
sont lus comme preuves historiques ; 9C3 fournit l'état de reproduction ultérieur.}
